\chapter{Literature Review}
\label{ch:literature-review}

The identification of a flying object from a measured trajectory is an inverse
problem: position and motion are observed, whereas target class and underlying
dynamics must be inferred. Traditional kinematic approaches use position,
velocity, acceleration, maneuvering behavior, or summary quantities extracted
from an estimated track. This thesis considers a complementary representation:
the time history of specific mechanical energy. The following review examines
the research underlying trajectory-based identification, physical energy
representations, reconstruction of degraded measurements, and machine-learning
classification. Its purpose is to motivate a testable research question rather
than assume that an energy representation is superior.

\section{Classification from Kinematic Observations}
\label{sec:lit-kinematic}

Motion-based target classification is closely connected to target tracking,
but the objectives are distinct. Tracking estimates a time-dependent state;
classification assigns a comparatively persistent identity or category.
Ristic, Gordon, and Bessell \cite{ristic2004kinematic} developed a framework
for classification using kinematic information, illustrating how
class-dependent motion characteristics can inform inference even without
a distinctive target signature. Angelova and Mihaylova
\cite{angelova2006joint} considered joint target tracking and classification
using kinematic radar measurements and probabilistic filtering. This line
of work makes measurement uncertainty and uncertainty about the dynamic
model central to the classification problem.

Auburn research provides a particularly relevant line of work on
classification from simulated projectile and missile trajectories.
Carpenter, Speakman, and Hartfield \cite{carpenter2016rapid}
studied rapid characterization of munitions using neural networks.
Zhou \cite{zhou2018deep} subsequently investigated deep learning
on simulated six-degree-of-freedom projectile histories, reporting
successful classification of three related short-range rocket classes
from early observations. Carpenter and colleagues
\cite{carpenter2019statistical} also examined statistical learning
for munition trajectory prediction, linking classification research
to the broader problem of extracting useful information from tracks.

A directly relevant study outside this research group is the dynamic
classification of ballistic missiles by Singh, Padmanabhan, and
Agarwal \cite{singh2014dynamic}.
Their use of neural networks and hidden Markov models demonstrates that
temporal patterns in observed motion can be used to distinguish missile
classes. However, it does not establish that a mechanical-energy history
is the best possible input representation. Jochumsen and colleagues
\cite{jochumsen2016recursive} also investigated recursive
random-forest classification from two-dimensional radar tracks,
illustrating an alternative connection between state estimation
and learning from kinematic observations.

Eckert, Carpenter, Hartfield, and Cervantes
\cite{eckert2020classification} extended early-flight classification
to simulated intermediate-range missile classes. Cervantes,
Moore, Carpenter, and Hartfield \cite{cervantes2024telemetry}
later compared classification from the first 10, 20, or 30 seconds
of telemetry against classification using summary quantities
such as trajectory maxima. Their analysis also considered feature
importance and missing data. This is a close methodological
precedent for comparing a compact physical representation
with direct time-history inputs, although their study did not
evaluate the specific energy history considered here.

These studies motivate examining whether a physically constructed
trajectory representation can make class-specific characteristics
more apparent to a classifier. They also caution against equating
good tracking accuracy with good class discrimination: the two tasks
must be evaluated separately.

\section{Mechanical Energy and Flight Dynamics}
\label{sec:lit-energy}

Classical orbital mechanics expresses the specific mechanical energy
of an object as the sum of kinetic and gravitational potential energy
per unit mass \cite{curtis2021orbital}. With gravitational potential
referenced to the Earth's surface, a convenient energy history is
\begin{equation}
    \mathcal{E}(t)
    = \frac{1}{2}v(t)^2
    + \mu\left(
        \frac{1}{R_E} - \frac{1}{R_E+h(t)}
      \right),
    \label{eq:lit-specific-energy}
\end{equation}
where $v$ is speed in the selected reference frame, $h$ is altitude,
$R_E$ is the reference Earth radius, and $\mu$ is the gravitational
parameter of Earth. The reference offset in potential energy does
not change energy differences.

In the ideal unpowered two-body problem, specific mechanical energy
computed using \emph{inertial} speed is constant. During atmospheric
flight, aerodynamic work reduces mechanical energy. Similarly, if
the speed is ground-relative rather than inertial, Earth's rotation
affects interpretation of changes in the calculated quantity.
A ground-relative energy history is therefore a physically motivated
\emph{feature}, not automatically a conserved invariant. This distinction
is particularly important when examining post-burnout behavior.

Powered flight can produce a characteristic increase in mechanical
energy, while the subsequent flight history may reflect both
nonconservative losses and the chosen reference frame. The evolution
of $\mathcal{E}(t)$ therefore offers a possible means of summarizing
changes in altitude and speed on a single physical scale. It may also
preserve information that an isolated apogee, maximum velocity,
or flight duration omits. Conversely, different three-dimensional
trajectories can have similar scalar energy histories; energy is not
a lossless replacement for the original trajectory. Its value as a
classification input must be determined experimentally.

An adjacent study by Evans, Carpenter, and Hartfield
\cite{evans2026clustering} explored unsupervised grouping of
solid-propellant missile designs using performance parameters
and maximum trajectory quantities. That study supports considering
derived physical summaries as representations of simulation output,
but its unsupervised objective and scalar inputs differ from
classification using an ordered energy history. It therefore
provides a relevant point of comparison rather than evidence
that energy profiles already improve supervised classification.

\section{Radar Measurements, Noise, and Incomplete Trajectories}
\label{sec:lit-noise}

A classifier operates on imperfect observations rather than exact
dynamics. Bar-Shalom, Li, and Kirubarajan
\cite{barshalom2001estimation} describe the estimation of dynamic
states from noisy observations, building on the state-space
filtering framework associated with Kalman \cite{kalman1960filtering}.
These methods clarify that uncertainty in position estimates
propagates into velocity and other derived quantities.

Three forms of degradation should be distinguished: measurement
perturbations, missing observations, and isolated gross errors.
Additive noise changes recorded values; sparse or intermittent
sampling removes segments of a trajectory; and outliers can
dominate an otherwise suitable regression model. These
degradations have different implications for reconstruction
and classification. In particular, numerical differentiation
can amplify high-frequency position error, and the kinetic
term in Eq.~\eqref{eq:lit-specific-energy} depends quadratically
on estimated speed.

Measurement degradation has also been addressed directly in
related missile-classification work. Eckert's subsequent thesis
\cite{eckert2020thesis} incorporated radar-equation-based
simulated measurement noise, an $\alpha$--$\beta$ filter,
and a comparison of neural networks with SVMs.
Cervantes and colleagues \cite{cervantes2024telemetry}
evaluated classification when telemetry values were removed
and imputed. Dove, Hartfield, and Carpenter
\cite{dove2025novelty} examined reduced observation windows
and missing data while distinguishing known missile classes
from previously unseen examples. These results demonstrate
the importance of testing missing-data sensitivity rather than
relying exclusively on intact simulation histories.

The present study employs repeatable synthetic degradation
experiments to investigate sensitivity to these effects.
Their purpose is a controlled robustness comparison. Such
experiments are not, without independent sensor validation,
proof that any selected synthetic error distribution describes
all operational radar systems.

\section{Trajectory Reconstruction and Robust Differentiation}
\label{sec:lit-reconstruction}

A continuous energy history generally requires a reconstructed
position and velocity history. State estimation offers one
approach \cite{kalman1960filtering,barshalom2001estimation};
smooth-function fitting offers another. Wahba
\cite{wahba1978smoothing} addressed spline smoothing
and the tradeoff between agreement with observations and
protection against model error. This motivates the use of
smooth representations when observations are irregular or sparse.

Ordinary least squares can be adversely influenced by a small
number of grossly erroneous samples. Rousseeuw and Leroy
\cite{rousseeuw1987robust} provide a foundation for robust
regression and outlier diagnosis. A physically constrained,
robust trajectory fit is consequently a plausible option
when measured tracks contain missing samples and isolated
large errors.

Derivatives require additional attention. Savitzky and Golay
\cite{savitzky1964smoothing} introduced local polynomial
smoothing and differentiation procedures. Their method
illustrates how derivatives can be estimated without relying
entirely on differences between adjacent noisy samples,
although the smoothing window may also suppress genuine
short-duration changes.

Gaubatz, Carpenter, and Hartfield
\cite{gaubatz2025fda} also considered functional-data-analysis
preprocessing for neural-network missile classification and
waypoint prediction. Treating sampled trajectories as
functions offers an alternative to feeding unprocessed
coordinate vectors directly into a classifier. Collins,
Carpenter, and Hartfield \cite{collins2025crud}
investigated coordinate-driven, robust, and uncertainty-aware
dynamic missile classification. These contributions
motivate comparisons among direct coordinates,
smoothed functional representations, and derived
physics-based quantities without treating preprocessing
as independent of the eventual classification objective.

For this thesis, reconstruction quality should not be assessed
solely through agreement in altitude or position. Velocity
error enters specific energy quadratically. It is therefore
important to assess the fidelity of derived energy profiles,
event timing, and classification performance, including
cases in which reconstruction fails.

\section{Supervised Classification Models}
\label{sec:lit-classification}

Support-vector machines (SVMs) are a natural baseline for
fixed-length energy-profile vectors. Cortes and Vapnik
\cite{cortes1995svm} developed the margin-based formulation
and its extension to nonlinear discrimination through
kernel functions. In a trajectory study, both feature scaling
and kernel hyperparameter tuning affect the resulting
decision boundary. Comparisons are meaningful only when
such decisions are made using training data rather than
the held-out evaluation set. In missile-specific work,
Eckert \cite{eckert2020thesis} compared SVMs and
neural networks using degraded simulated flight histories,
finding that a radial-kernel SVM performed particularly
well on the investigated dataset. Dove, Hartfield,
and Carpenter \cite{dove2025novelty} combined
one-class SVMs with linear discriminant analysis to
address the different challenge of identifying both
known classes and unfamiliar trajectories. Neither
result alone establishes the best classifier for an
energy-based representation.

Breiman's random forests \cite{breiman2001forests} provide
an alternative using ensembles of randomized trees.
They are useful for evaluating compact sets of physically
motivated features, including event times and energy
statistics, and have also been explored for radar-track
classification \cite{jochumsen2016recursive}. They
do not establish that an apparently predictive feature
represents a causal physical distinction. In related
solid-rocket-motor modeling, Sheils and colleagues
\cite{sheils2025explainability} used SHAP analyses to
interpret the influence of model inputs on predicted
motor performance. Their work motivates examining
which energy-profile regions or derived descriptors
drive classification decisions, rather than reporting
accuracy alone.

Neural networks can instead learn nonlinear combinations
of profile samples. LeCun, Bengio, and Hinton
\cite{lecun2015deep} describe the broader approach of
representation learning. A multilayer perceptron can
ingest a resampled energy vector directly, whereas a
convolutional or recurrent architecture can exploit
temporal structure. Long short-term memory networks
\cite{hochreiter1997lstm} are one example of a model
designed to capture dependencies across sequential data.
The distinction between architecture and input
representation is particularly important in light of
DiMaggio, Hartfield, and Carpenter's
\cite{dimaggio2026ensemble} study of genetically
optimized neural-network ensembles for classifying
similar solid rocket motor variants from early
flight histories. Their results show that
network architecture and aggregation can affect
classification performance when classes are similar;
they do not isolate the effect of an energy representation.
These are possible comparisons, not an a priori
argument that deeper models must perform better.

\section{Time-Series Representation and Evaluation}
\label{sec:lit-timeseries}

Energy histories are time series, and their encoding changes
the question a classifier is asked to solve. The early-flight
windows examined by Cervantes and colleagues
\cite{cervantes2024telemetry}, the functional
preprocessing considered by Gaubatz and colleagues
\cite{gaubatz2025fda}, and the early-flight
ensembles of DiMaggio and colleagues
\cite{dimaggio2026ensemble} provide
particularly relevant precedents for comparing
time-history encodings on common observation windows. Bagnall and
colleagues \cite{bagnall2017bakeoff} emphasize consistent
experimental comparison of time-series methods, while
Ismail Fawaz and colleagues \cite{fawaz2019review}
review deep-learning models for time-series classification.

Sampling profiles on a common absolute-time grid retains
information about when events occur, but requires an explicit
policy for unequal trajectory lengths. Normalizing each
profile to a common duration allows shape-based comparison
but can remove differences in event timing or total flight
duration. Restricting observations to powered flight or
apogee changes the available evidence and may make
early-stage classification possible, at the cost of
excluding later behavior.

A central experimental safeguard is to separate
\emph{underlying trajectories}, not merely their noisy
realizations. Multiple degraded versions of one simulation
should not appear on both sides of a training/test split,
because doing so can exaggerate generalization to new
trajectories. Robustness comparisons should report
classification metrics, per-class behavior, and the rate
of unsuccessful reconstructions rather than scoring
successful cases alone.

\section{Research Gap and Thesis Direction}
\label{sec:lit-gap}

Prior work establishes that kinematic motion histories can
be used for classification \cite{ristic2004kinematic,
singh2014dynamic,jochumsen2016recursive,
eckert2020classification,cervantes2024telemetry},
that uncertain tracks require appropriate reconstruction
\cite{barshalom2001estimation,rousseeuw1987robust,
eckert2020thesis,collins2025crud},
and that temporal representations can be learned or
classified using multiple model families
\cite{bagnall2017bakeoff,fawaz2019review,
gaubatz2025fda,dimaggio2026ensemble}.
Earlier work with physical performance summaries and
unsupervised grouping \cite{evans2026clustering},
as well as methods for detecting classes absent from
training \cite{dove2025novelty}, provides further
context for defining the scope of the present study.
These findings motivate, but do not independently validate,
the energy-based formulation proposed in this thesis.

The research question is whether a specific-mechanical-energy
representation offers a favorable balance of class
discrimination, resistance to measurement degradation,
data requirements, and computational cost relative
to direct kinematic inputs. The subsequent methodology
will compare reference and reconstructed histories,
standalone profiles and hybrid descriptors, and
classical classifiers and neural networks. Those
comparisons will determine the merits and limits
of the representation for the available datasets.
